\scanpage{77}
\mspage{14}
du groupe structural de $\unsure{\mathcal G_3|\mathfrak S_3}$ à $\unsure{\mathcal G_3|\mathfrak S_3}^+\simeq\mathbb Z/3\mathbb Z$.

Quand \deleted{toutes} les conditions précédentes sont vérifiées, y compris la donnée d'un \inserted{tel} coloriage, alors on trouve un rev. \inserted{principal} canonique \deleted{$\unsure{X'^+\subset X'}$} \inserted{$X'^+\subset X'$} de $X$, $(*)$ de groupe $\mathbb Z/3\mathbb Z$ -- dont la classe d'isomorphie est un élément
\[
(\unsure{43|45})\qquad \operatorname{cl}(X)\in H^1(X,\mathbb Z/3\mathbb Z)
\]
qu'on peut considérer comme l'obstruction à l'existence d'une pondération, pour la carte triangulée $X$ : sa nullité est néc. et suffisante pour cette existence. De plus, les pondérations pour lesquelles les faces $+$ soient les faces directes (pour l'orientation donnée) sont en corr. $1$-$1$ avec les sections du \deleted{fibré} rev. principal précédent de groupe $\mathbb Z/3\mathbb Z$.
\figure{\begin{tikzpicture}\draw (0,0)--(.75,-1.1);\end{tikzpicture}}
\margin{76}
